\documentclass[exercice,a4]{thedocument}%
\usepackage{texmaths-tikz}%
\usepackage{pstricks}%
\usepackage{pst-bezier}%
\usepackage{pstricks-add}%
\usepackage{pst-eucl}%
\usepackage{pst-func}%
\usepackage{pst-tree}%
\usepackage{pst-circ}%
\usepackage{tikz}%
\usepackage{pgfplots}%
\usepackage{tkz-tab}%
\usepackage{tikzpagenodes}%
\usepackage{tabularray}%
\usepackage{tabularx}%
\usepackage{multirow}%
\usepackage{multicol}%
\usepackage{scratch3}%
\usepackage{enumitem}%
\usepackage[np]{numprint}%
\usepackage{pas-tableur}%
\usepackage{listings}%
\usepackage{graphicx}%
\usepackage{adjustbox}%
\usepackage{needspace}%
\usepackage{xcolor}%
\usepackage{eurosym}%
\usepackage{diagbox}%
\usepackage{colortbl}%
\usepackage{multido}%
\usepackage{icomma}%
\usepackage[french]{babel}%
\startdatakeys
    \source{26GENMATG11}
    \BO{2026}
    \cycle{4}
    \competencesDuSocle{CA,RA,CO}
    \objectifsApprentissage{C4A1aa,C4A2ad,C4B2cb,C4D2aa,C4D3ac}
    \motsCles{Pythagore,droites parallèles,Thalès,périmètre,construction}
\enddatakeys
\definecolor{gris}{gray}{0.5}%
\providecommand{\wideparen}[1]{\overset{\frown}{#1}}%
\begin{document}%
\providecommand{\degres}{^\circ}%
Dans cet exercice, on considère la figure ci-contre.

	\medskip

	\begin{minipage}{0.47\linewidth}
Les points A, B et M sont alignés.

Les points A, C et N sont alignés.

Le triangle ABC est rectangle en B.

Le triangle AMN est rectangle en M.
\end{minipage}\hfill
%%début PSTricks
%\begin{minipage}{0.49\linewidth}
%\psset{unit=1cm,arrowsize=2pt 3}
%\begin{pspicture}(7,4.8)
%	%\psgrid
%	\psline(0,0.6)(7,4.7)
%	\psline(1,1.2)(6,1.2)(6,4.1)%AMN
%	\psline(4,1.2)(4,2.96)%BC
%	\uput[dl](1,1.2){A} \uput[r](6,1.2){M} \uput[dr](6,4.1){N}
%	\uput[ur](4,1.2){B} \uput[ur](3.9,3.13){C}
%	\psline[linewidth=0.8pt]{<->}(1,0.8)(4,0.8)\uput[d](2.5,0.8){6,4 cm}
%	\psline[linewidth=0.8pt]{<->}(4,0.8)(6,0.8)\uput[d](5,0.8){3,2 cm}
%	\psline[linewidth=0.8pt]{<->}(0.9,1.3)(3.95,3.11)\rput{31}(2.5,2.45){8 cm}
%	\psframe(4,1.2)(3.75,1.45)
%	\psline(6,1.2)(6,0)%M..
%\end{pspicture}
%\end{minipage}
%%Fin PSTricks
%Tikz
\begin{tikzpicture}[baseline={(current bounding box.center)},line join=round, line cap=round]
\coordinate (A) at (0,0);
\coordinate (B) at (4,0);
\coordinate (M) at (5.5,0);
\coordinate (C) at (4,2.4);
\coordinate (N) at (5.5,3.3);
\draw[line width=1.2pt,red,fill=red!30] (4,0) rectangle ++(-0.4,0.4)
 	(5.5,0) rectangle ++(-0.4,0.4) ;
\draw[line width=1.2pt]
	(A) -- (M)
	(B) -- (C)
	(-1,-0.6) -- (6.5,3.9)
	(5.5,-0.8) -- (5.5,3.9);
\foreach \coor/\nom/\pos in {(A)/A/below, (B)/B/{above right},
	(C)/C/above, (M)/M/{below right}, (N)/N/{below right}}
	{\draw[line width=1pt, shift={\coor}] (2pt,2pt)--(-2pt,-2pt) (2pt,-2pt)--(-2pt,2pt) (0,0) node[\pos]{\nom} ;}
	\draw[<->, line width=0.8pt, > = stealth,shift={(0,-0.5)}] (0,0) -- (4,0) node[pos=0.5,below,sloped]{6,4 cm} ;
	\draw[<->, line width=0.8pt, > = stealth,shift={(0,-0.5)}] (5.5,0) -- (4,0) node[pos=0.5,below,sloped]{3,2 cm} ;
	\draw[<->, line width=0.8pt, > = stealth,shift={(120:0.5)}] (0,0) -- (4,2.4) node[pos=0.5,above,sloped]{8 cm} ;
\end{tikzpicture}\hfill~
%Fin tikz


\medskip

On donne :

AB $= 6,4$ cm ; BM $= 3,2$ cm et AC $= 8$ cm.

\emph{La figure n'est pas en vraie grandeur.}

\begin{enumerate}
\item Tracer sur la copie le triangle ABC en vraie grandeur et en laissant les traits de construction.
\item Démontrer que BC $= 4,8$ cm.
\item Justifier que les droites (BC) et (MN) sont parallèles.
\item Démontrer que MN $= 7,2$ cm et AN $= 12$ cm.
\item Le périmètre du triangle ABC et le périmètre du quadrilatère BMNC sont-ils égaux ?\\
\textbf{Argumenter la réponse en précisant la démarche.}
\end{enumerate}

\needspace{10\baselineskip}
\medskip
\showthesource
\end{document}
